% ============================================================================
% WORKBOOK — Section 6.4: Constructing Triangles from Three Measurements
% CCSS 7.G.A.2
% Practice-focused reinforcement for the study guide topic
% ============================================================================
\section{Constructing Triangles from Three Measurements}
\topicTitle{Constructing Triangles from Three Measurements}

\begin{quickReview}{Triangle Conditions}
Three measurements can produce a \textbf{unique} triangle, \textbf{many} triangles, or \textbf{none}:

\begin{center}
\begin{tabular}{l l}
    \textbf{Given} & \textbf{Result} \\
    \hline
    Three sides (SSS) & Unique (if triangle inequality holds) \\
    Two sides + included angle (SAS) & Unique \\
    Two angles + included side (ASA) & Unique \\
    Two angles + non-included side (AAS) & Unique \\
    Three angles (AAA) & Many (same shape, different sizes) \\
\end{tabular}
\end{center}

\textbf{Key facts:}
\begin{itemize}[leftmargin=*]
    \item Triangle inequality: each side $<$ sum of the other two.
    \item Angle sum: the three angles of a triangle always add to $180°$.
\end{itemize}
\end{quickReview}

% ── Warm-Up ──────────────────────────────────────────────────────────────────
\begin{practiceBox}[funGreen]{\faSun~Warm-Up}

\practiceHeader[funGreen]{Find the Missing Angle}
The three angles of a triangle sum to $180°$. Find the missing angle.

\begin{multicols}{2}
\prob Angles: $50°$, $60°$, \answerBlank[2cm]
\answerExplain{$70°$}{$180° - 50° - 60° = 70°$. The three angles must sum to $180°$.}

\prob Angles: $90°$, $45°$, \answerBlank[2cm]
\answerExplain{$45°$}{$180° - 90° - 45° = 45°$. This is an isosceles right triangle.}

\prob Angles: $30°$, $30°$, \answerBlank[2cm]
\answerExplain{$120°$}{$180° - 30° - 30° = 120°$. The triangle has two equal angles and one obtuse angle.}

\prob Angles: $72°$, $54°$, \answerBlank[2cm]
\answerExplain{$54°$}{$180° - 72° - 54° = 54°$. Two angles are equal, making it isosceles.}

\prob Angles: $60°$, $60°$, \answerBlank[2cm]
\answerExplain{$60°$}{$180° - 60° - 60° = 60°$. All three angles are equal — equilateral triangle.}
\end{multicols}
\end{practiceBox}

% ── Practice A: Possible or Impossible? ──────────────────────────────────────
\begin{practiceBox}{Can This Triangle Exist?}

\practiceHeader{Write Yes or No. If yes, state whether the triangle is unique, or if many are possible.}

\prob Angles $40°$, $60°$, $80°$. \answerBlank[3cm]
\answerExplain{Yes, many triangles}{$40 + 60 + 80 = 180°$ \checkmark. But three angles (AAA) allow many triangles of different sizes with the same shape.}

\prob Angles $50°$, $60°$, $80°$. \answerBlank[3cm]
\answerExplain{No}{$50 + 60 + 80 = 190°$. The angles sum to more than $180°$, which is impossible for a triangle.}

\prob Sides $8$, $15$, $6$. \answerBlank[3cm]
\answerExplain{No}{Check: $8 + 6 = 14$, but $14 < 15$. The triangle inequality fails because the two shorter sides cannot reach around the longest side.}

\prob Sides $8$, $15$, $10$. \answerBlank[3cm]
\answerExplain{Yes, unique}{Check: $8 + 10 = 18 > 15$ \checkmark, $8 + 15 = 23 > 10$ \checkmark, $10 + 15 = 25 > 8$ \checkmark. Three valid sides (SSS) give exactly one triangle.}

\prob Two sides $7$ and $11$, included angle $65°$. \answerBlank[3cm]
\answerExplain{Yes, unique}{Two sides and the included angle (SAS) always produce exactly one triangle.}

\prob Angles $90°$, $100°$, $-10°$. \answerBlank[3cm]
\answerExplain{No}{An angle cannot be negative. Every angle of a triangle must be greater than $0°$. Also, $90 + 100 = 190° > 180°$.}
\end{practiceBox}

% ── Practice B: Classify the Condition ───────────────────────────────────────
\begin{practiceBox}[funTeal]{Name the Condition}

\practiceHeader[funTeal]{Identify each set of given information as SSS, SAS, ASA, AAS, or AAA.}

\begin{multicols}{2}
\prob Sides: $5$, $7$, $9$. \answerBlank[2cm]
\answerExplain{SSS}{All three sides are given — this is Side-Side-Side (SSS).}

\prob Sides: $6$, $10$; included angle: $50°$. \answerBlank[2cm]
\answerExplain{SAS}{Two sides and the angle between them — this is Side-Angle-Side (SAS).}

\prob Angles: $40°$, $70°$; included side: $8$. \answerBlank[2cm]
\answerExplain{ASA}{Two angles and the side between them — this is Angle-Side-Angle (ASA).}

\prob Angles: $35°$, $55°$, $90°$. \answerBlank[2cm]
\answerExplain{AAA}{All three angles are given — this is Angle-Angle-Angle (AAA).}

\prob Angles: $30°$, $80°$; non-included side: $12$. \answerBlank[2cm]
\answerExplain{AAS}{Two angles and a side not between them — this is Angle-Angle-Side (AAS).}
\end{multicols}
\end{practiceBox}

% ── Practice C: True or False ────────────────────────────────────────────────
\begin{practiceBox}[funPurple]{True or False?}

\trueOrFalse{If three angles of a triangle add up to $180°$, you can draw exactly one triangle.}
\answerExplain{False}{The angles sum correctly, but AAA does not determine the size. You can draw infinitely many triangles with those same angles but different side lengths.}

\trueOrFalse{SSS (three sides) always produces a unique triangle.}
\answerExplain{False}{SSS produces a unique triangle \emph{only if} the triangle inequality holds. If the sides don't satisfy the inequality, no triangle exists at all.}

\trueOrFalse{If two angles of a triangle are $45°$ and $65°$, the third angle is $70°$.}
\answerExplain{True}{$180° - 45° - 65° = 70°$. There is exactly one value for the third angle.}

\trueOrFalse{SAS (two sides and the included angle) always determines a unique triangle.}
\answerExplain{True}{Two sides and the angle between them fix all remaining measurements. Only one triangle can be drawn.}
\end{practiceBox}

% ── Word Problems ────────────────────────────────────────────────────────────
\begin{practiceBox}[funBlue]{\faGlobe~Word Problems}

\wordProblem{An engineer knows two beams are $12$ ft and $16$ ft, and the angle between them must be $90°$. How many different triangle shapes are possible?}{triangles}
\answerExplain{Exactly $1$}{Two sides and the included angle (SAS) always produce one unique triangle. This is a right triangle with legs $12$ and $16$.}

\wordProblem{A surveyor measures three angles of a triangular plot: $48°$, $62°$, and $70°$. Can this be a real triangle?}{Yes or No}
\answerExplain{Yes}{$48 + 62 + 70 = 180°$ \checkmark. The angles sum to exactly $180°$, so triangles with these angles exist (but AAA means many sizes are possible).}

\wordProblem{Mia says she can draw a triangle with sides $3$ cm, $4$ cm, and $8$ cm. Is she correct? Why or why not?}{}
\answerExplain{No}{$3 + 4 = 7$, and $7 < 8$. The triangle inequality fails because the two shorter sides cannot reach around the longest side.}
\end{practiceBox}

% ── Challenge ────────────────────────────────────────────────────────────────
\begin{challengeBox}
\prob Two angles of a triangle are in the ratio $2:3$. The third angle is $50°$. Find all three angles, and explain whether this gives a unique triangle or many.
\answerExplain{$52°$, $78°$, $50°$; many triangles}{The first two angles sum to $180° - 50° = 130°$. In ratio $2:3$: $2x + 3x = 130$, so $5x = 130$ and $x = 26$. Angles: $52°$, $78°$, $50°$. Since we only know angles (AAA), many triangles of different sizes are possible.}

\prob Can you have a triangle where all three sides are different and all three angles are $60°$? Explain.
\answerExplain{No}{If all three angles are $60°$, the triangle is equilateral, which means all three sides must be equal. It is impossible to have three different sides with three equal angles.}
\end{challengeBox}

\encouragement{You are a triangle construction expert! Keep building your geometry skills!}
